\chapter{Successor Extensions and Compiled Workflows}\label{chap:extensions}

This chapter catalogs successor-side capabilities that are compiled from
shared vehicle contracts into deterministic runtime studies and artifacts. It
is intentionally separate from the reconstructed 1995 TAOS syntax. A feature
listed here is a TAORYX extension unless its source, implementation, and
verification boundary are explicitly identified as historical.

\section{Compiled Capability Catalog}\label{sec:compiled-capability-catalog}

The catalog below is the index for major compiled capabilities. Operational
details belong in the corresponding Markdown architecture and plan pages; this
chapter records the physical meaning, compiled surface, and evidence required
before a result is promoted.

\begin{longtable}{p{0.20\linewidth}p{0.25\linewidth}p{0.42\linewidth}}
\toprule
\textbf{Capability} & \textbf{Compiled surface} & \textbf{Evidence and boundary}\\
\midrule
\endhead
Reachability envelope &
  Search commands, terminal criteria, fidelity tier, and deterministic batch
  propagation. &
  JSON result contains the complete search space, successful and unsuccessful
  samples, timeout classifications, terminal margins, and replay metadata.\\
Deployment and ejection &
  Optional stage mass, separation event, impulse frame, residual policy,
  detached-body definition, and accepted-boundary child spawn. &
  Parent pre/post state, child initial state, event identity, mass accounting,
  impulse, and momentum residual are recorded when this capability is
  composed with a study. Invalid requests are explicit deployment failures.\\
Ballistic body families &
  Sphere, cylinder, spheroid, and cone geometry with projected area, inertia,
  center of pressure, and tumbling policy. &
  Geometry and coefficient assumptions are serialized. Surrogate aerodynamics
  are not treated as validated test or CFD data.\\
Fidelity ladder &
  Point-mass 3-DOF, pseudo-6-DOF, and native rigid-body 6-DOF child tiers. &
  Tier comparisons retain the same search coordinates and identify which
  parent and child models actually ran.\\
Runtime model collection &
  Stable model IDs, parent IDs, deferred spawn requests, and atomic insertion
  into the active model collection. &
  Accepted event boundaries are the only deployment points; rejected solver
  trials do not create models or mutate the collection.\\
Successor language and lowering &
  Explicit dynamics modes, runtime declarations, deployment requests, and
  source-located extension diagnostics. &
  Historical and successor grammar profiles remain distinct; unsupported
  features fail closed instead of becoming implicit physics.\\
Native aero and table models &
  Force and moment coefficient tables, direct dimensional wrench tables,
  atmosphere providers, and Earth-rotation frame adapters. &
  Table domains, source quality, margins, and frame conventions are retained
  with the run artifact.\\
Guidance, control, and thermal gates &
  Attitude moments, actuator saturation, LQR design, terminal guidance,
  heating limits, and objective scoring. &
  Controller results are downstream of plant validation; gates and quality
  scores remain separate claims.\\
Vehicle onboarding and validation &
  Registry contracts, source anchors, trim points, parity cases, and fidelity
  ladder promotion. &
  A new vehicle must pass source, geometry, dynamics, and artifact gates before
  it is treated as an integrated example.\\
Result and visualization compilation &
  JSON, CSV telemetry, Matplotlib trajectories, coverage, capability,
  event-timeline, projected-area, and fidelity-comparison artifacts. &
  A manifest joins source configuration, resolved parameters, search space,
  event IDs, telemetry, plots, and terminal classifications.\\
\bottomrule
\end{longtable}

\section{Reachability Envelopes}\label{sec:compiled-reachability}

A reachability envelope is a standalone, finite, explicitly parameterized search
rather than a claim that every point in continuous control space has been tested.
It does not require staging, deployment, ejection, or spawned bodies. An
air-breathing aircraft, drone, glider, or spacecraft can produce a reachability
envelope with no deployment feature configured. Each
candidate retains its decision vector, search coordinates, terminal state,
success status, failure reasons, path metrics, and witness trajectory when
requested. The result format also distinguishes impact, timeout, invalid
integration, and constraint failure.

Terminal acceptance may include impact radius and impact-speed limits. A
candidate that reaches the horizon before ground contact is a timeout, not an
impact. Timeout candidates can be extracted and rerun with a longer horizon
without reconstructing the original search manually.

The canonical X-15 exercise composes the two features: Reachability supplies
the search and terminal evaluation, while Deployment supplies the booster
release and spawned-body histories. It runs one shared launch grid through
three declared tiers: point-mass 3-DOF, pseudo-6-DOF, and a
reduced-parent/native-rigid-child 6-DOF deployment tier. The X-15 result is an
X-15-scaled surrogate integration case, not evidence of a native historically
validated X-15 batch model.

\section{Deployment and Ejection}\label{sec:compiled-deployment}

Deployment and ejection are a separate accepted-boundary capability. They can
be used by the general runtime without a reachability study, and can be
optionally composed with Reachability when parent/child outcomes belong in the
same envelope. Deployment is an accepted-boundary event. The parent is propagated to the event
time, its pre-event state is retained, mass and declared impulse discontinuities
are applied, and a child is initialized from the same accepted timestamp. The
parent and child then continue as independent models with shared provenance.
The generic child is a full runtime vehicle: it may be propulsive, guided,
instrumented, aerodynamic, ballistic, or a spacecraft. A
\texttt{DetachedBodyDefinition} is only the optional specification for an
aero-ballistic child, not a requirement for deployment.
For Alpha 3, the supported and documented authoring path is the passive
aero-ballistic specialization; higher-level configuration of arbitrary active
children is deferred to Alpha 4.

The shared contract distinguishes liquid, solid, hybrid, and unknown
propulsion behavior. Throttle and cutoff requests are rejected when the
declared hardware cannot realize them. Separation mechanism, impulse frame,
ejection energy, detached mass, inertia, center of mass, center of pressure,
and body shape are explicit rather than inferred from a vehicle-specific
adapter.

A body-frame impulse is transformed through the parent attitude. The retained
and detached bodies receive equal-and-opposite translational impulses under
the declared partition. The event artifact records the parent pre/post
velocity, child initial velocity, and momentum residual so that an apparently
successful trajectory cannot hide an invented momentum jump.

\section{Tumbling and Rigid-Body Promotion}\label{sec:compiled-tumbling}

The pseudo-6-DOF tier is appropriate for fast screening and for exposing
orientation-dependent projected-area effects. The native rigid-body child tier
uses quaternion attitude, body rates, diagonal principal inertia, aerodynamic
force, center-of-pressure moments, and explicit unresolved-rate damping. The
damping is a stability-bounded ballistic surrogate; it is not a replacement
for a measured six-degree-of-freedom coefficient database.

The mathematical basis for the projected-area and rotational closures is in
\cref{sec:ballistic-tumbling}. The implementation research and generated
coefficient studies are maintained under \texttt{analysis/tumbling}; the
runtime contract is documented under \texttt{docs/architecture} and the Alpha
3 deployment plan.

\section{Artifacts and Promotion Gates}\label{sec:compiled-artifacts}

Every promoted envelope must emit a machine-readable manifest and the
corresponding evidence set:

\begin{itemize}
  \item parent and child trajectories with deployment markers;
  \item accepted-event timeline with mass, impulse, mechanism, and identity;
  \item successful and unsuccessful capability footprint;
  \item search-space coverage and uncovered candidate bins;
  \item projected-area and aerodynamic-load histories;
  \item fidelity comparison for matched candidates; and
  \item parent and per-child telemetry tables with terminal classification.
\end{itemize}

Promotion requires deterministic event-boundary behavior, valid quaternion
norm, bounded timestep refinement, mass and momentum closure, explicit timeout
handling, and a claim boundary that identifies surrogate assumptions. A green
plot is not sufficient evidence when the search space, parameterization, or
failed candidates are absent.
